\section{基于岩石识别的避障策略}
\label{sec:ch5_obstacle_avoidance}

岩石识别的最终目标不是“看见障碍”，而是“安全通过障碍区”。
因此，本节将第\ref{subsec:ch5_reconstruction_counting}节输出的岩石几何信息与第4章动力学风险量统一到可执行的避障决策中，形成“识别结果 $\rightarrow$ 风险评估 $\rightarrow$ 控制输出”的闭环策略。

\subsection{避障决策逻辑}
\label{subsec:ch5_decision_logic}

综合沉陷、碰撞、能耗和时延四类风险，定义局部决策代价为
\begin{equation}
\mathcal{R}_{total}(x,y)=0.35R_{sink}(x,y)+0.30R_{col}(x,y)+0.20R_{eng}(x,y)+0.15R_{delay}(x,y),
\label{eq:ch5_total_risk}
\end{equation}
其中权重取自表5-2的参数设置。对应风险触发规则如下：
\begin{myitemize}
	\item 当 $R_{sink}>0.6$ 时，降低速度并提高局部重规划频率；
	\item 当 $R_{col}>0.5$ 时，禁止直行策略，仅允许绕行或等待；
	\item 当 $\mathrm{SOC}<20\%$ 时，提高能耗项惩罚权重；
	\item 当链路时延超过阈值时，启用保守安全距离模式。
\end{myitemize}

在候选策略集 $\Pi=\{\pi_{left},\pi_{right},\pi_{wait}\}$ 上，选择最小风险策略
\begin{equation}
\pi^\star=\arg\min_{\pi\in\Pi}\left[\sum_{k=0}^{H-1}\mathcal{R}_{total}\!\left(\mathbf{x}_k^\pi\right)+\lambda_l L(\pi)+\lambda_u\Delta u(\pi)\right],
\label{eq:ch5_policy_select}
\end{equation}
其中 $H$ 为预测步长，$L$ 为路径代价，$\Delta u$ 为控制变化惩罚。

\begin{table}[H]
	\centering
	\small
	\caption{数字孪生预演下候选避障策略风险评分}
	\label{tab:ch5_candidate_policy_risk}
	\begin{tabular}{p{0.30\textwidth}p{0.20\textwidth}}
		\toprule
		候选策略 & 风险评分 \\
		\midrule
		左绕路径 $\pi_{left}$ & $0.741215$ \\
		右绕路径 $\pi_{right}$ & $0.785823$ \\
		等待路径 $\pi_{wait}$ & $0.789683$ \\
		\bottomrule
	\end{tabular}
\end{table}

由表\ref{tab:ch5_candidate_policy_risk}可见，在当前地形与障碍配置下左绕策略风险最低，因此被选为执行策略。

\subsection{路径调整算法}
\label{subsec:ch5_path_adjustment}

\begin{figure}[H]
	\centering
	\includegraphics[width=0.90\textwidth]{ch5/ch5_fig_5_3_obstacle_rehearsal.png}
	\vspace{0.4em}
	
	\begin{minipage}{0.90\textwidth}
		\footnotesize
		\textbf{图例注记：}
		\fcolorbox{black!30}{cyan!35}{\rule{1.2em}{0.8em}}~低风险通行区；
		\fcolorbox{black!30}{yellow!45}{\rule{1.2em}{0.8em}}~中风险过渡区；
		\fcolorbox{black!30}{red!45}{\rule{1.2em}{0.8em}}~高风险障碍影响区；
		\textcolor{blue!70!black}{\rule{1.2em}{1pt}}~左绕候选路径；
		\textcolor{orange!85!black}{\rule{1.2em}{1pt}}~右绕候选路径；
		\textcolor{gray!70!black}{\rule{1.2em}{1pt}}~等待策略轨迹。
	\end{minipage}
	\caption{数字孪生环境下候选避障策略预演结果}
	\label{fig:ch5_obstacle_rehearsal}
\end{figure}

图\ref{fig:ch5_obstacle_rehearsal}给出了统一起终点与障碍分布条件下的候选策略预演结果。
图中颜色由冷到暖对应累计风险由低到高，路径叠加结果显示左绕策略主要穿越低风险通行走廊，
而右绕与等待策略分别受到高坡度区域与时延惩罚项影响，导致综合代价上升。
该图与表\ref{tab:ch5_candidate_policy_risk}中的定量评分相互印证，说明“风险场评估-策略筛选”机制能够在执行前有效识别更优规避方向。

局部路径以全局路径为先验，在障碍区附近执行代价重加权与增量平滑更新：
\begin{equation}
\mathbf{p}_{i}^{t+1}=\mathbf{p}_{i}^{t}
+\eta_1\left(\mathbf{p}_{i}^{0}-\mathbf{p}_{i}^{t}\right)
+\eta_2\left(\mathbf{p}_{i-1}^{t}+\mathbf{p}_{i+1}^{t}-2\mathbf{p}_{i}^{t}\right)
-\eta_3\nabla \mathcal{R}_{total}(\mathbf{p}_{i}^{t}),
\label{eq:ch5_local_path_update}
\end{equation}
其中前三项分别对应原路径保持、曲率平滑与风险场排斥。

\begin{algorithm}[H]
\caption{基于岩石识别的局部避障决策}
\label{alg:ch5_obstacle_avoidance}
\begin{algorithmic}[1]
\REQUIRE 全局路径 $\mathcal{P}_g$，岩石集合 $\mathcal{O}$，风险场 $\mathcal{R}_{total}$
\ENSURE 局部控制命令 $\mathbf{u}_k$
\STATE 从 $\mathcal{P}_g$ 截取前视窗口，生成候选策略集 $\Pi$
\FOR{$\pi\in\Pi$}
\STATE 计算轨迹滚动风险与路径代价，得到 $J(\pi)$
\ENDFOR
\STATE 选择 $\pi^\star=\arg\min_{\pi\in\Pi}J(\pi)$
\IF{$J(\pi^\star)>\tau_{risk}$}
\STATE 触发局部重规划并降低速度上限
\ENDIF
\STATE 按式\eqref{eq:ch5_local_path_update}更新局部轨迹
\STATE 输出控制命令 $\mathbf{u}_k$
\end{algorithmic}
\end{algorithm}

\subsection{安全距离计算模型}
\label{subsec:ch5_safe_distance_model}

在式\eqref{eq:ch5_safe_distance_base}基础上，进一步考虑滑移与沉陷导致的制动距离放大，得到动态安全距离模型：
\begin{equation}
d_{safe}=r_{rover}+r_{rock}+v\tau+k_\sigma \sigma_p+k_s z_{mean}+k_\mu s_{mean},
\label{eq:ch5_safe_distance_dynamic}
\end{equation}
其中 $z_{mean}$ 与 $s_{mean}$ 分别为局部窗口内平均沉陷与平均滑移率。

该模型使安全边界可随动力学状态实时扩张或收缩。
当地面承载能力降低（$z_{mean}$ 增大）或驱动稳定性下降（$s_{mean}$ 增大）时，系统自动增大避障裕度，避免“几何可通行但动力学不可通行”的失效场景。
